\section{Compile-time ORM архитектура}
\label{sec:architecture}

Compile-time ORM архитектурата на библиотеката е организирана около ясно разграничение между логическия модел на данните, междинното представяне на заявките и backend-специфичния слой за материализация. Именно този пласт носи първата основна научна ос на разработката: как компилаторът може да участва в моделирането, ограничаването и валидирането на достъпа до данни още преди runtime фазата. За да бъде постигнато това, архитектурата е изградена от няколко слоя, всеки със собствена отговорност и със строго определени интерфейси.

В теоретичната глава беше показано, че унифицирането е възможно на равнището на логическия модел и query IR-а, но не и чрез игнориране на различията между backend-ите. След корутинната теория от предходната глава тук фокусът съзнателно се връща към compile-time ядрото на системата: entity описанията, query IR-а, \code{connector\_trait<DB>} и capability механизма. По този начин библиотеката се разглежда не просто като идея, а като организирана система от типове, contracts и формални граници на допустимите операции.

\subsection{Архитектурни цели и ограничения}

Архитектурата е проектирана при едновременно действие на няколко ограничения. Първо, user-facing API трябва да остане в рамките на идиоматичен C++, без въвеждане на отделен DSL и без runtime reflection механизъм. Второ, query структурата трябва да бъде представима като \code{constexpr} обект, така че значима част от валидирането да се извършва от компилатора. Трето, backend слойът трябва да бъде разширяем без наследяване от общ виртуален интерфейс, тъй като подобен интерфейс би скрил различията между backend-ите и би изместил проверките към runtime.

Четвъртото ограничение е свързано със самия multi-backend характер на системата. Библиотеката трябва да може да поддържа SQL и non-SQL конектори без фалшиво уеднаквяване. Това означава, че архитектурата трябва едновременно да допуска общ API и да изразява ограниченията на конкретен backend по формален и проверим начин. Оттук естествено следва изборът на trait-базиран connector слой и capability механизъм.

\begin{figure}[H]
\centering
\begin{tikzpicture}[node distance=1.3cm and 1.4cm]
    \node[ormaccent] (app) {\textbf{Приложен слой}\\entity типове и query композиция};
    \node[ormblock, below=of app] (entity) {\textbf{Entity слой}\\\code{property}, \code{relationship}, \code{table\_name\_trait}};
    \node[ormblock, below=of entity] (ir) {\textbf{Query IR слой}\\\code{field}, \code{Rule}, \code{select\_query}, \code{insert\_query}};
    \node[ormblock, below=of ir] (connector) {\textbf{Connector слой}\\\code{connector\_trait<DB>} и capability tags};
    \node[ormblock, below=of connector] (driver) {\textbf{Driver / runtime слой}\\SQL/BSON/Redis/Cypher/CQL материализация};

    \draw[ormline] (app) -- (entity);
    \draw[ormline] (entity) -- (ir);
    \draw[ormline] (ir) -- (connector);
    \draw[ormline] (connector) -- (driver);
\end{tikzpicture}
\caption{Слоеста архитектура на библиотеката}
\label{fig:architecture_layers}
\end{figure}

Фигура~\ref{fig:architecture_layers} показва, че библиотеката следва последователна вертикална организация. Приложният код не комуникира директно с драйвера; той работи с C++ типове и query builder-и. Драйверът от своя страна не познава домейн семантиката, а само materialized формата, генерирана от connector слоя. Именно поради това архитектурата може да остане едновременно строга и надграждаема.

\subsection{Слоеста организация и разделяне на отговорностите}

Най-горният слой е приложният код, който дефинира entity типове, изгражда заявки чрез \code{select}, \code{insert}, \code{update} и \code{deleteq}, и ги подава на \code{db<DB>}. Следващият слой е логическият модел на библиотеката --- entity описанията и query IR. Третият слой е connector abstraction-ът, където \code{connector\_trait<DB>} материализира същата логика към backend-специфичен синтаксис и поведение. Най-долу стои конкретният драйвер, wire protocol или mock implementation.

Тази архитектура има важен практически ефект. Ако потребителят промени \code{SQLiteDB} с \code{MongoDB}, query IR остава логически същият, но конекторът го тълкува като filter/projection изпълнение вместо SQL рендиране. Това не означава, че всеки IR е допустим върху всеки backend; означава, че \emph{формата на абстракцията} е обща, а допустимостта се определя формално чрез capability механизма.

\begin{table}[H]
\centering
\small
\caption{Основни архитектурни слоеве и техните отговорности}
\label{tab:architecture_responsibilities}
\begin{tabularx}{\textwidth}{L{2.8cm}L{4.0cm}Y}
\toprule
\textbf{Слой} & \textbf{Основни артефакти} & \textbf{Отговорност} \\
\midrule
Приложен & entity типове, query builder повиквания & описва домейн логика и формулира операции в идиоматичен C++ \\
Логически & \code{property}, \code{relationship}, \code{field}, \code{Rule} & моделира данните и заявките независимо от конкретен backend \\
Connector & \code{connector\_trait<DB>}, capability tags & проверява допустимостта и материализира IR-а към backend поведение \\
Runtime/driver & конкретни connection handle-и и драйверни API-та & изпълнява материализираната заявка и връща резултати \\
\bottomrule
\end{tabularx}
\end{table}

Таблица~\ref{tab:architecture_responsibilities} показва, че библиотеката не използва ``голям обект'', който да прави всичко. Вместо това различните решения са локализирани в отделни слоеве. Това е важно както за надграждаемостта, така и за академичния анализ, защото позволява отделните твърдения да бъдат проследени до конкретни типове и модули.

\subsection{Entity слой като архитектурна основа}

Entity слоят е изграден върху \code{property<T, Name>}, \code{relationship<...>} и \code{table\_name\_trait<T>}. Той има двойна роля. От една страна, представя C++ домейн модела. От друга страна, служи като логическо описание на схемата или структурата на персистентните данни. Това означава, че архитектурата не прави отделна runtime registry от метаданни. Вместо това типовата система съдържа достатъчно информация, за да бъдат извлечени имена на полета, типове, връзки и целеви контейнер за съхранение.

Особено важен е фактът, че entity слоят е независим от вида база данни. Нито \code{property}, нито \code{relationship} са SQL-специфични. Те описват логически атрибути и отношения. Така се създава основата за сравнимо поведение на релационни и нерелационни backend-и. На архитектурно ниво именно това прави възможно библиотеката да твърди, че C++ моделът е първичният носител на логическата схема.

\subsection{Query IR като централен архитектурен компонент}

Query IR е централният архитектурен компонент. Той обединява field референции, предикати, placeholder-и и допълнителни клаузи като \code{order\_by}, \code{group\_by} и \code{limit}. Вместо да се работи с текст на заявка, библиотеката изгражда типизиран обект, чийто шаблонни параметри кодират логическата структура на операцията.

От архитектурна гледна точка това решение има няколко последствия. Първо, валидирането на структурата се извършва рано. Второ, query IR може да бъде анализиран от connector слоя без загуба на семантична информация. Трето, подготвените заявки могат да запазят не SQL низ, а самото междинно представяне, което е по-естествено за библиотека, ориентирана към C++ типовата система.

Четвърто, query IR-ът е единствената зона, в която се срещат entity описанието и по-късната backend материализация. В този смисъл той е архитектурният ``шарнир'' на цялата библиотека: достатъчно висок, за да остане storage-agnostic, и достатъчно конкретен, за да може да бъде изпълнен.

\subsection{Конекторен слой и формалният contract}

Формално connector слоят е реализиран чрез специализация на \code{connector\_trait<DB>}. Този подход е аналогичен на стандартни trait шаблони като \code{std::iterator\_traits}. Типът \code{DB} може да бъде прост wrapper над connection handle, а цялата ORM логика за него да остане в trait специализацията. По този начин библиотеката не изисква потребителят или авторът на конектора да наследява общ базов клас.

\code{connector\_trait<DB>} определя поне три ключови аспекта: \code{wire\_type}, \code{cursor\_type} и \code{execute(...)}. По желание специализацията декларира и capability tags като \code{supports\_joins}, \code{supports\_transactions}, \code{supports\_aggregation} или \code{supports\_concurrent\_execute}. Наличието или отсъствието на тези псевдо-маркери е достатъчно, за да може по-горният слой да наложи compile-time ограничения.

\begin{lstlisting}[language=C++,caption={Извадка от формалния connector contract},label={lst:connector_contract}]
template <typename DB>
concept is_connector = requires {
    typename connector_trait<DB>::template wire_type<int>::type;
    typename connector_trait<DB>::cursor_type;
};

template <typename DB, typename Cap>
inline constexpr bool has_capability =
    detail::capability_check<DB, Cap>::value;
\end{lstlisting}

Listing~\ref{lst:connector_contract} показва, че contract-ът е структурен, а не наследствен. Това е решаващ архитектурен избор. Ако съществуваше виртуален интерфейс с runtime dispatch, част от типово проверимата информация щеше да бъде изгубена. Вместо това библиотеката предпочита compile-time структура, при която изискванията са видими и за компилатора, и за автора на конектора.

\begin{figure}[H]
\centering
\begin{tikzpicture}[node distance=1.35cm and 1.25cm]
    \node[ormaccent] (query) {\textbf{Типизирана заявка}\\join, transaction, aggregation intent};
    \node[ormblock, below=of query] (db) {\code{db<DB>} фасада};
    \node[ormblock, below left=of db] (capok) {Capability налице\\заявката е допустима};
    \node[ormdanger, below right=of db] (capfail) {Capability липсва\\compile-time отказ};
    \node[ormblock, below=of capok] (exec) {\code{connector\_trait<DB>::execute}};

    \draw[ormline] (query) -- (db);
    \draw[ormline] (db) -- (capok);
    \draw[ormline] (db) -- (capfail);
    \draw[ormline] (capok) -- (exec);
\end{tikzpicture}
\caption{Capability gating между query IR-а и connector слоя}
\label{fig:capability_gating}
\end{figure}

Фигура~\ref{fig:capability_gating} е важна, защото показва къде архитектурата взема решение за допустимост. Това не е драйверът и не е runtime логът. Решението се взема в публичния API слой, който вече има достъп до query IR-а и до trait информацията за конектора.

\begin{table}[H]
\centering
\small
\caption{Роля на основните capability tags в архитектурата}
\label{tab:capability_tags}
\begin{tabularx}{\textwidth}{L{3.6cm}Y}
\toprule
\textbf{Capability} & \textbf{Архитектурен смисъл} \\
\midrule
\code{supports\_joins} & разрешава join-базираните query форми за съответния backend \\
\code{supports\_transactions} & позволява използване на \code{transaction\_guard} и \code{begin}/\code{commit}/\code{rollback} hooks \\
\code{supports\_aggregation} & обозначава поддръжка на aggregation операции в логическия query слой \\
\code{supports\_embedding} & прави експлицитна native поддръжка на вградени структури \\
\code{supports\_upsert} & маркира backend-и, при които upsert е част от естествения execution contract \\
\code{supports\_bulk\_insert} & обозначава наличието на по-ефективен път за масово записване \\
\bottomrule
\end{tabularx}
\end{table}

\subsection{Потребителският фасаден слой db<DB>}

Класът \code{db<DB>} е публичният фасаден слой. Той капсулира препратка към конкретния backend обект и предоставя уеднаквен API за изпълнение. На практика този слой има две роли. Първата е удобство за потребителя: библиотеката се използва чрез един и същ интерфейс независимо от backend-а. Втората е архитектурен филтър: точно тук се извършват compile-time проверки за capability и съвместимост.

Следователно \code{db<DB>} не е просто тънък wrapper над \code{execute}. Той е мястото, където библиотеката формализира границата между разрешено и неразрешено поведение за избрания backend. Точно този слой позволява съобщения за грешка от вида ``конекторът не поддържа \code{JOIN} операции'', вместо грешката да се появи по-късно като неясно runtime поведение.

\begin{lstlisting}[language=C++,caption={Извадка от фасадния слой \code{db<DB>}},label={lst:db_facade}]
template <typename DB>
    requires is_connector<DB>
class db
{
public:
    template <typename Query>
    auto operator<<(Query q)
    {
        if constexpr (is_select_query<Query>)
        {
            if constexpr (decltype(q.join_clauses())::size > 0)
            {
                static_assert(has_capability<DB, cap::supports_joins>,
                    "orm::db: this connector does not support JOIN operations.");
            }
        }
        return connector_trait<DB>::execute(*conn_, std::move(q));
    }
};
\end{lstlisting}

Listing~\ref{lst:db_facade} показва защо \code{db<DB>} е повече от удобна фасада. Той е формалното място, където query структурата и connector capability информацията се срещат. Именно тук compile-time архитектурата се превръща в реално ограничение върху API употребата.

\subsection{Миграции като архитектурно разширение}

Архитектурата включва и прототипен migration слой чрез \code{migrate<DB>}, \code{live\_schema}, \code{entity\_meta} и DDL операции като \code{create\_table\_op}, \code{add\_column\_op}, \code{drop\_column\_op} и \code{alter\_column\_type\_op}. Този слой е особено важен от гледна точка на тезата, защото прави явна връзката между C++ модела и реалната схема на съхранение.

Миграционният слой не е универсално завършен за всеки backend, но архитектурно показва, че библиотеката не се ограничава само до query execution. При достатъчно богата специализация на \code{connector\_trait<DB>} entity описанието може да участва и в еволюцията на схемата. Това променя и самия смисъл на ORM слоя: той вече не е само интерфейс към данните, а потенциален източник на schema-aware tooling.

\subsection{Реализация на основните подсистеми}

След като архитектурният модел е очертан, следващият въпрос е как той се материализира в конкретни C++ конструкции и как отделните подсистеми си взаимодействат в хранилището. Настоящият раздел възстановява липсващата инженерна перспектива: не само какви слоеве съществуват, а как точно се натрупват \code{where} клаузите, как query обектът се запазва за повторна употреба, как \code{migrate<DB>} превръща разлика между entity описанието и \code{live\_schema} в DDL операции и как helper слоят налага правила за нишкова безопасност и транзакции.

Тук се фиксира и работната нотация, без да се връща отделна речникова глава. В настоящия раздел \code{DB} означава generic backend тип, \code{db<DB>} обозначава публичната фасада за изпълнение, \code{connector\_trait<DB>} е формалният конекторен contract, \code{field<...>} е compile-time представяне на поле, \code{Placeholder<T>} е анонимен runtime параметър, а \code{orm::ph<T, \_N>} е индексиран placeholder. Репозиторийният контекст на тази механика е концентриран главно в \path{lib/include/ORM/entity/}, \path{lib/include/ORM/query/}, \path{lib/include/ORM/connector/}, \path{lib/include/ORM/db/migration/} и във верификационния слой под \path{tests/unit/} и \path{tests/integration/}.

\subsubsection{Мини казус с типовете User и Post}

За да остане изложението непосредствено обвързано с реалния код, този раздел използва като опорна mini case study типовете \code{User} и \code{Post} от \path{tests/integration/test_mockdb.cpp}. Те са особено полезни, защото съчетават скаларни полета, явно именуване на логическите контейнери и join сценарий. Същият пример по-късно се използва и за placeholder binding, prepared query поведение и SQL rendering assertions.

\begin{lstlisting}[language=C++,caption={Entity типове от \protect\path{tests/integration/test_mockdb.cpp}},label={lst:mockdb_entities}]
struct Post
{
    orm::property<int, "id">             id;
    orm::property<int, "author_id">      author_id;
    orm::property<std::u8string, "body"> body;
};

struct User
{
    orm::property<int, "id">            id;
    orm::property<std::u8string, "name"> name;
    orm::property<double, "score">      score;
};
\end{lstlisting}

Listing~\ref{lst:mockdb_entities} показва не само как се описва entity моделът, но и как той непосредствено участва в тестовете. \code{table\_name\_trait<User>} и \code{table\_name\_trait<Post>} фиксират логическите контейнери \code{"users"} и \code{"posts"}, а member pointer-и като \code{\&User::id} и \code{\&Post::author\_id} по-късно участват в query builder-а. Това е пряка илюстрация на един от централните мотиви на дипломната работа: моделът, query формата и execution тестовете използват едни и същи compile-time артефакти.

\subsubsection{Скаларни свойства и именуване на контейнерите}

Реализацията на \code{property<T, Name>} свързва стойностен тип с compile-time име. Функцията \code{column\_name()} дава еднозначен достъп до символното име на полето и се използва от конекторите при рендиране към SQL колони, document fields, key-space сегменти или други backend-специфични понятия. Конструкцията е минималистична, но достатъчна, защото не изисква външна runtime регистрация на метаданни.

Именуването на логическия контейнер за съхранение се реализира чрез \code{table\_name\_trait<T>}. Макар името да е наследено от ORM традицията, архитектурно то играе по-обща роля: в SQLite, PostgreSQL и MySQL означава таблица, в MongoDB --- collection, в Redis --- key prefix, а в Neo4j --- label. Така библиотеката използва единен логически интерфейс за адресиране на физически различни структури.

Практическата последица е важна: авторът на конектора не извлича схема от външен registry, а от типовата система. Това намалява boilerplate кода и прави грешките в именуването по-лесни за локализиране както в компилационната фаза, така и в unit тестовете.

\subsubsection{Field wrappers, правила и placeholder-и}

Полето в query израз се представя чрез \code{field<\&T::m>}. Това е \code{constexpr} wrapper над member pointer, който носи достатъчно информация за типа на контейнера, типа на стойността и името на колоната. Операторите върху \code{field} изграждат типизирани \code{Rule} обекти. Така изрази от вида \code{field<\&User::id> == Placeholder<int>\{\}} не са низове, а compile-time описания на предикати.

Placeholder механизмът има две форми. \code{Placeholder<T>} моделира анонимни параметри, които се консумират позиционно от аргументите на \code{execute()}. \code{orm::ph<T, std::placeholders::\_N>} моделира индексирани параметри, които позволяват една runtime стойност да се използва на повече от една позиция. Това е особено важно за backend-и като SQLite, PostgreSQL и MySQL, при които placeholder слотът е отделна част от materialization логиката.

Тестовете \code{IndexedPlaceholderRendersQuestionMarkN}, \code{IndexedPlaceholderReusedSameArgInSQL} и \code{TwoDistinctIndexedPlaceholdersRenderCorrectSlots} в \path{tests/integration/test_mockdb.cpp} показват, че тази подсистема не е теоретична добавка, а реално използван механизъм за по-точни parameter binding сценарии.

\subsubsection{Изграждане на query обекти}

Фабричните функции \code{select}, \code{insert}, \code{update} и \code{deleteq} създават query обекти, върху които последователно се наслагват \code{where}, \code{join}, \code{order\_by}, \code{group\_by} и \code{limit}. Същественото е, че всяка от тези операции не променя низ, а връща нов типизиран обект с разширено вътрешно състояние. В резултат query builder-ът остава едновременно fluent на вид и статичен по природа.

Това позволява в тестовете да се проверява не само крайният резултат от рендирането, а и самата форма на междинното представяне --- например броят на \code{where} клаузите, типът на response tuple-а или коректността на \code{GROUP BY} и \code{ORDER BY} composition-а. В \path{tests/integration/test_mockdb.cpp} това се вижда в \code{SelectWithInnerJoin}, \code{SelectWithOrderByDesc}, \code{GroupByMultipleFields}, \code{SelectWithLimit} и множество други сценарии, които валидират различни фази на query composition.

Ключов имплементационен детайл е, че builder веригата е \code{constexpr}-съвместима. Вътрешният помощник \code{orm::detail::tuple\_cat} използва разгъване на пакет чрез индексни последователности вместо зависимост от runtime конкатенация. Практическата последица е, че значима част от query формата може да бъде изградена и проверена още по време на компилация.

\begin{lstlisting}[language=C++,caption={Query обект деклариран като compile-time константа},label={lst:constexpr_query}]
using namespace orm::literals;
static constexpr auto get_active =
    orm::select(orm::field<&User::id>, orm::field<&User::name>)
        .where(orm::field<&User::score> > 0.0)
        .where(orm::field<&User::id> == orm::Placeholder<int>{})
        .limit(10_per_page & 2_page);
\end{lstlisting}

Операторът \code{\&} комбинира помощниците \code{\_per\_page} и \code{\_page} в \code{Pagification} обект, чиято структура е известна по време на компилация. Хедърът \path{lib/include/ORM/query/limits.hpp} поддържа и симетрична форма чрез \code{*}, но използването на \code{\&} в Listing~\ref{lst:constexpr_query} е по-четливо за заявъчния стил на библиотеката. Така заявката може да бъде декларирана като \code{static constexpr} и да се използва многократно без реконструиране на builder формата при всяко извикване.

\begin{figure}[H]
\centering
\begin{tikzpicture}[node distance=1.35cm and 1.2cm]
    \node[ormaccent] (api) {\textbf{Fluent API}\\\code{select}/\code{where}/\code{join}/\code{limit}};
    \node[ormblock, below=of api] (typed) {Ново типизирано\\query състояние при всяка стъпка};
    \node[ormblock, below left=of typed] (shape) {Compile-time\\структура на заявката};
    \node[ormblock, below right=of typed] (params) {Runtime\\стойности на параметрите};
    \node[ormblock, below=of $(shape)!0.5!(params)$] (exec) {\code{connector\_trait<DB>::execute}};

    \draw[ormline] (api) -- (typed);
    \draw[ormline] (typed) -- (shape);
    \draw[ormline] (typed) -- (params);
    \draw[ormline] (shape) -- (exec);
    \draw[ormline] (params) -- (exec);
\end{tikzpicture}
\caption{Взаимодействие между query composition, compile-time структура и runtime параметри}
\label{fig:query_composition_execution}
\end{figure}

Фигура~\ref{fig:query_composition_execution} показва, че реализацията работи в два режима едновременно: структурната форма на заявката е фиксирана в типа, докато конкретните стойности се подават отделно към \code{execute(...)}. Именно това прави възможно едновременното наличие на compile-time проверки и runtime параметризация.

\subsubsection{Изпълнение, резултати и достъп до полета}

На етап изпълнение \code{db<DB>} предава query IR-а към \code{connector\_trait<DB>::execute}. Резултатът се връща като \code{orm::result<Row, FieldTuple>}. В тази конструкция \code{Row} е материализираният tuple тип, а \code{FieldTuple} запазва projection metadata-то за избраните полета. Това решение позволява както итериране върху материализираните редове, така и достъп чрез \code{get\_field<\&T::m>} или позиционен достъп чрез \code{get<I>}. Реализацията така остава type-safe и след изпълнението на заявката.

Тестовете \code{SelectResultRowTypeIsTuple}, \code{SelectMultiFieldRowType}, \code{SelectGetByMemberPointer} и \code{SelectGetByIndex} показват, че result слоят не е просто контейнер от анонимни редове, а структурирана абстракция с типово проследима проекция. Това е важно за тезата, защото демонстрира, че типовата дисциплина не приключва с query builder-а, а продължава и в result handling-а.

\subsubsection{Prepared queries като повторно използваем execution артефакт}

\code{prepared\_query<DB, Query>} разширява изпълнителния модел. Вместо да кешира SQL низ, той съхранява връзка към backend обекта и самия query IR. Това е естествено продължение на compile-time подхода: библиотеката третира query обекта като първичен артефакт, а не като временна стъпка по пътя към текстов резултат.

\begin{lstlisting}[language=C++,caption={Извадка от \code{prepared\_query<DB, Query>}},label={lst:prepared_query}]
template <typename DB, typename Query>
class prepared_query
{
public:
    prepared_query(DB& conn, Query q) : conn_(conn), query_(std::move(q)) {}

    template <typename... Params>
    auto execute(Params&&... params) const
    {
        return connector_trait<DB>::execute(
            conn_, query_, std::forward<Params>(params)...);
    }

    [[nodiscard]] const Query& query() const noexcept { return query_; }
};
\end{lstlisting}

Listing~\ref{lst:prepared_query} показва важен инженерен избор: повторната употреба е организирана около IR-а, а не около текстовата материализация. Това позволява една и съща логическа заявка да бъде изпълнявана многократно с различни параметри, без да се реконструира builder формата. Точно това поведение се валидира от \code{PrepareReturnsExecutableObject}, \code{PrepareWithParamExecutesCorrectly}, \code{PrepareCanBeCalledMultipleTimesWithDifferentParams} и \code{PrepareExposesQueryIR}.

\begin{figure}[H]
\centering
\begin{tikzpicture}[node distance=1.35cm and 1.2cm]
    \node[ormaccent] (build) {Създаване на query IR};
    \node[ormblock, right=of build] (prepare) {\code{db.prepare(query)}};
    \node[ormblock, right=of prepare] (store) {Съхранен IR +\\референция към \code{DB}};
    \node[ormblock, below=of store] (exec1) {\code{execute(1)}};
    \node[ormblock, left=of exec1] (exec2) {\code{execute(99)}};
    \node[ormblock, right=of exec1] (exec3) {\code{execute(7)}};

    \draw[ormline] (build) -- (prepare);
    \draw[ormline] (prepare) -- (store);
    \draw[ormline] (store) -- (exec1);
    \draw[ormline] (store) -- (exec2);
    \draw[ormline] (store) -- (exec3);
\end{tikzpicture}
\caption{Lifecycle на \code{prepared\_query}: един IR, множество изпълнения}
\label{fig:prepared_query_lifecycle}
\end{figure}

\subsubsection{Миграционен слой и DDL генериране}

Реализацията на \code{migrate<DB>} сравнява описаната от entity типовете схема с \code{live\_schema}. Когато има разлика, се създават DDL операции и се подават към \code{connector\_trait<DB>::ddl\_for(...)}. Това е добър пример за clean separation на responsibilities: diff логиката е обща, а конкретният DDL синтаксис остава backend-специфичен.

\begin{lstlisting}[language=C++,caption={Ядро на migration diff pipeline-а},label={lst:migration_diff}]
[[nodiscard]] std::vector<ddl_op> diff(
    const live_schema& live,
    std::span<const entity_meta> entities) const
{
    std::vector<ddl_op> ops;

    for (const auto& entity : entities)
    {
        auto live_it = live.find(entity.table_name);
        if (live_it == live.end())
        {
            ops.push_back(create_table_op{entity.table_name, entity.columns});
            continue;
        }

        for (const auto& [col, type] : entity.columns)
        {
            if (live_it->second.find(col) == live_it->second.end())
                ops.push_back(add_column_op{entity.table_name, col, type});
        }
    }
    return ops;
}
\end{lstlisting}

Тук ясно се вижда разликата между \emph{минимален работещ конектор} и \emph{пълен connector contract}. Един backend може да изпълнява заявки, без да поддържа DDL генериране. Но ако трябва да участва в schema-aware tooling, неговият contract трябва да бъде разширен с \code{ddl\_for(...)} hooks и последователна migration верификация.

\begin{figure}[H]
\centering
\begin{tikzpicture}[node distance=1.35cm and 1.2cm]
    \node[ormaccent] (entity) {Entity metadata\\\code{entity\_meta}};
    \node[ormblock, right=of entity] (live) {Runtime schema\\\code{live\_schema}};
    \node[ormblock, below=of $(entity)!0.5!(live)$] (diff) {\code{migrate<DB>::diff}};
    \node[ormblock, below left=of diff] (ops) {DDL операции\\\code{create/add/drop/alter}};
    \node[ormblock, below right=of diff] (ddl) {\code{generate\_ddl}};
    \node[ormblock, below=of $(ops)!0.5!(ddl)$] (run) {dry-run или execution path};

    \draw[ormline] (entity) -- (diff);
    \draw[ormline] (live) -- (diff);
    \draw[ormline] (diff) -- (ops);
    \draw[ormline] (diff) -- (ddl);
    \draw[ormline] (ops) -- (run);
    \draw[ormline] (ddl) -- (run);
\end{tikzpicture}
\caption{Поток на schema migration подсистемата}
\label{fig:migration_pipeline}
\end{figure}

Тестовете в \path{tests/unit/test_schema_migration.cpp} покриват създаване на таблица, добавяне и отпадане на колони, dry-run кодове на завършване, генериране на DDL и state-machine поведението на \code{run()}. Това е силно доказателство, че migration слоят не е само концепция, а реално проследим pipeline.

\subsubsection{Нишкова безопасност и transaction helpers}

В библиотеката съществува отделен слой за \code{connection\_pool<DB, N>}, \code{thread\_local\_db<DB>} и \code{transaction\_guard<DB>}. Тези компоненти показват, че реализацията не спира до еднократно изпълнение на заявка, а разглежда и operational сценарии с повторна употреба на връзки, изолиране по нишки и rollback дисциплина. Важното архитектурно решение е, че достъпът до тези механизми също е capability-базиран. Например \code{connection\_pool} изисква \code{supports\_concurrent\_execute}, а \code{transaction\_guard} --- \code{supports\_transactions}.

\begin{lstlisting}[language=C++,caption={Извадка от thread-safety helper слоя},label={lst:thread_safety_helpers}]
template <typename DB, std::size_t N>
class connection_pool
{
    static_assert(
        requires { typename connector_trait<DB>::supports_concurrent_execute; },
        "connection_pool requires "
        "connector_trait<DB>::supports_concurrent_execute. "
        "Add 'using supports_concurrent_execute = void;' "
        "to connector_trait<DB>.");

public:
    [[nodiscard]] connection_guard<DB> acquire();
};

template <typename DB>
class transaction_guard
{
    static_assert(
        requires { typename connector_trait<DB>::supports_transactions; },
        "transaction_guard requires "
        "connector_trait<DB>::supports_transactions. "
        "Add 'using supports_transactions = void;' "
        "to connector_trait<DB>.");
};
\end{lstlisting}

Така compile-time contract-ът не се ограничава до query формата, а достига и до lifecycle политиките на конектора. Това е важно за аргумента на тезата, че библиотеката не е просто builder за заявки, а инфраструктура за систематична и безопасна работа с backend ресурси.

\begin{table}[H]
\centering
\small
\caption{Основни подсистеми и тяхната тестова опора}
\label{tab:core_subsystem_tests}
\begin{tabularx}{\textwidth}{L{3.1cm}L{4.1cm}Y}
\toprule
\textbf{Подсистема} & \textbf{Основни файлове} & \textbf{Тестова опора} \\
\midrule
Entity моделиране & \path{ORM/entity/property.hpp}, \path{ORM/entity/relationship.hpp} & \path{tests/unit/test_property.cpp}, \path{tests/unit/test_relationship.cpp} \\
Query composition & \path{ORM/query/select.hpp} и свързаните query headers & \path{tests/unit/test_query.cpp}, \path{tests/integration/test_mockdb.cpp} \\
Prepared queries & \path{ORM/connector/prepared_query.hpp} & \code{Prepare*} сценариите в \path{tests/integration/test_mockdb.cpp} \\
Schema migration & \path{ORM/db/migration/migration.hpp} & \path{tests/unit/test_schema_migration.cpp} \\
Thread safety & \path{ORM/db/connectors/ThreadSafety/thread_safety.hpp} & \path{tests/unit/test_thread_safety.cpp} \\
\bottomrule
\end{tabularx}
\end{table}

\subsubsection{Извод за реализацията}

Основните подсистеми имат пряка опора в кода и тестовете. \path{tests/unit/test_property.cpp} и \path{tests/unit/test_relationship.cpp} валидират entity описанията. \path{tests/unit/test_query.cpp} доказва, че query builder-ът действително конструира правилни compile-time структури. \path{tests/integration/test_mockdb.cpp} и \path{tests/integration/test_sqlite.cpp} валидират изпълнението, prepared query сценариите и result handling-а. \path{tests/unit/test_schema_migration.cpp} покрива diff логиката на миграциите, а \path{tests/unit/test_thread_safety.cpp} проверява lifecycle слоя за конкурентен достъп.

Следователно реализацията на основните подсистеми не е само проектно намерение, а архитектурно решение, подкрепено от систематична верификация. Това е и най-важният извод от настоящия раздел: compile-time ORM архитектурата в хранилището не остава на равнище абстракция, а се материализира в работещи подсистеми с ясно разграничени отговорности и проследима репозиторийна опора.

\subsection{Архитектурни инварианти}

На базата на разгледаните слоеве могат да се формулират няколко инварианта. Първо, приложният код никога не бива да зависи пряко от драйверен API. Второ, query логиката се формулира в типизиран IR, а не в backend-native низове. Трето, разширяването към нов backend се извършва чрез trait специализация и capability декларации, а не чрез промяна на самия query API. Четвърто, ако една операция е недопустима за даден backend, това трябва да стане видимо възможно най-рано --- за предпочитане при компилация.

Тези инварианти правят архитектурата едновременно инженерно подредена и академично защитима. Те позволяват всяко следващо разширение да бъде оценявано не само по това дали ``работи'', а и по това дали запазва вече формулираните слоеве и граници.

\subsection{Архитектурен извод}

Обобщено, архитектурата на библиотеката разделя проблема на три равнища: логически модел, типизиран query IR и backend materialization. Това разделение позволява едновременно строга типова дисциплина, формализирана разширяемост и контролирано поведение върху различни модели за съхранение. В следващата глава фокусът се премества към втората архитектурна ос --- асинхронното изпълнение --- където става ясно как същото compile-time ядро се свързва с coroutine lifecycle, thread-pool offload и native non-blocking интеграция.
